% Fuente: docs/adr/0085-in-browser-evidence-validation.md (resultados por componente);
% crates/evidence-verify/src/report.rs y web/apps/qadra/src/lib/evidence/presentation.mjs.
\begin{tikzpicture}[
  comp/.style={caja, text width=3.75cm, minimum height=0.6cm, align=left},
  esc/.style={caja, text width=3.3cm, minimum height=0.95cm, align=left},
  r/.style={font=\diag, text=ESCOMGris}]

  % Componentes
  \draw[zona] (0,0) rectangle (4.45,-7.7);
  \node[zonaLabel] at (0,0) {COMPONENTES};
  \foreach \i/\t in {
      1/{Paquete},
      2/{Sello de tiempo del documento},
      3/{Firma del documento},
      4/{Vinculación y consentimiento\\{\diag no aplica al sello del servidor}},
      5/{Certificados y raíz de confianza},
      6/{Revocación},
      7/{Manifiesto y su sello de tiempo\\{\diag solo en la versión 2}},
      8/{Fechas y registros}} {
    \node[comp] (c\i) at (2.225,{-0.2-0.88*\i}) {\t};
    \draw[ESCOMAzulOscuro, line width=0.6pt] (c\i.east) -- (4.8,{-0.2-0.88*\i});
  }
  \draw[ESCOMAzulOscuro, line width=0.6pt] (4.8,-1.08) -- (4.8,-7.24);

  % Cinco resultados
  \node[cajaClara, text width=2.6cm, align=left, minimum height=3.3cm] (cinco) at (6.55,-4.16)
    {\textbf{Cada uno obtiene uno de cinco resultados}\\[3pt]
     Correcto\\ Fallido\\ Falta evidencia para decidirlo\\ No compatible\\ No aplica};
  \draw[flecha] (4.8,-4.16) -- (cinco.west);

  % Precedencia
  \node[r, anchor=south west, align=left] at (8.4,-0.5) {Resultado global: se toma el primero que aplica};
  \node[esc, cajaError] (f) at (10.15,-1.2) {\textbf{Fallido}: algún componente obligatorio falló};
  \node[esc] (nc) at (10.15,-2.7) {\textbf{No compatible}: un obligatorio no es compatible o ningún perfil coincide};
  \node[esc] (p) at (10.15,-4.2) {\textbf{Parcial}: falta evidencia en algún obligatorio};
  \node[esc, cajaOsc] (c) at (10.15,-5.7) {\textbf{Completo}: todos los obligatorios comprobados};
  \draw[flecha] (cinco.east) -- ++(0.2,0) |- (f.west);
  \draw[flecha] (f) -- node[etq, right=1pt]{si no} (nc);
  \draw[flecha] (nc) -- node[etq, right=1pt]{si no} (p);
  \draw[flecha] (p) -- node[etq, right=1pt]{si no} (c);

  % Informativo y aclaración
  \node[cajaGris, text width=2.4cm] (act) at (14.65,-1.2) {\textbf{Estado actual}\\ vigencia y revocación hoy, con el reloj del equipo};
  \draw[flechaSuave] (act.west) -- node[etq, below=3pt, align=center]{informativo:\\ no cambia\\ el resultado} (f.east);
  \node[caja, text width=2.4cm, align=left] (nota) at (14.65,-4.7)
    {\textbf{Firma correcta con perfil de certificado fallido}: el resultado no es completo, pero eso no significa por sí solo que el PDF se haya alterado};

  \node[r, text width=10.6cm, align=left, anchor=north west] at (5.25,-6.65)
    {Se separan tres preguntas: si la firma corresponde a los bytes exactos, si los certificados cumplen el perfil y encadenan a una raíz fija, y si falta evidencia para decidir.};
\end{tikzpicture}
